\documentclass[11pt,a4paper]{article}
\usepackage[utf8]{inputenc}
\usepackage[T1]{fontenc}
\usepackage{lmodern}
\usepackage[margin=2.4cm]{geometry}
\usepackage{amsmath,amssymb}
\usepackage{bm}
\usepackage{booktabs}
\usepackage{xcolor}
\usepackage[colorlinks=true,linkcolor=blue!55!black,urlcolor=blue!55!black]{hyperref}
\usepackage{caption}
\captionsetup{font=small,labelfont=bf}
\setlength{\parskip}{0.35em}
\newcommand{\code}[1]{\texttt{#1}}
\let\oldunderscore\_
\renewcommand{\_}{\oldunderscore\allowbreak}
\setlength{\emergencystretch}{2.5em}

\title{\vspace{-1.2cm}Counting every degree of freedom of the plasma of \code{TransitionListener} exactly once\\[0.2em]
\large Notes on the degrees of freedom, the transverse photon and the scale factor}
\date{September 2026}
\author{Carlo Tasillo}

\begin{document}
\maketitle
\vspace{-1.4cm}

\section{What is counted where}

\code{TL} keeps two inventories for the degrees of freedom in the plasma.
The fields appearing in the effective potential are counted mode by mode, each polarization of
a gauge boson separately, with their field- and temperature-dependent masses. Everything else is counted by the two
tabulated radiation baths, the kinetically coupled bath \code{kin\_coupled\_*}
(by default the Standard Model bath) and
a kinetically decoupled bath \code{kin\_decoupled\_*} (which can be a secluded dark sector).
Every particle which does not appear in the effective potential must appear in
exactly one of the two.
For a potential
that contains Standard Model fields (for instance for the 2HDM models these are $W$, $Z$, $t$, $b$, $\tau$, $h$ and
the Goldstone modes, which are flagged as \code{is\_SM}) the same fields were counted twice in
some places and not at all in others: the energy density counted the photon both in the
potential and in the tabulated bath, while the entropy counted essentially none of them: the
frozen tables still hold what is left of the $\tau$ and the $b$ at the freezing temperature,
but nothing of $W$, $Z$, $t$ and $h$.

In the Landau gauge the modes appearing in the potential are not necessarily physical particles:
each gauge boson
contributes two transverse and one longitudinal mode, each broken generator adds a Goldstone
mode, and each gauge boson comes with a ghost that removes one bosonic degree of freedom. For
the Standard Model fields of the 2HDM at the zero-temperature vacuum,
%
\begin{align}
  \underbrace{1}_{h} + \underbrace{1+2}_{G^0, G^\pm}
  + \underbrace{2+1}_{Z_\text{T}, Z_\text{L}}
  + \underbrace{2+1}_{\gamma_\text{T}, \gamma_\text{L}}
  + \underbrace{4+2}_{W_\text{T}, W_\text{L}}
  - \underbrace{4}_{\text{ghosts}} = 12
  = \underbrace{1}_{h} + \underbrace{3}_{Z} + \underbrace{6}_{W^\pm} + \underbrace{2}_{\gamma} ,
\end{align}
%
so the gauge-dependent modes cancel and the physical count remains; with the fermions $t$, $b$
and $\tau$ included, this sums to $12 + \tfrac{7}{8}\cdot 28 = 36.5$ degrees of freedom at
temperatures far above their masses.
That number is the bridge between the two inventories.

\section{The three corrections}

\paragraph{1. The modes of the potential.} \code{h\_eff\_DS} and \code{g\_eff\_DS}, i.e.~the
entropy and energy degrees of freedom, were originally calculated by masking out
every field flagged \code{is\_SM}, skipping the Goldstone modes and keeping the ghosts. They now
count every mode and subtract one ghost per gauge boson, which is the counting
\code{radiationEnergyDensity} has always used. For the 2HDM the masked fields were counted
nowhere: 68 entropy degrees of freedom at $T = 50\,\text{GeV}$ instead of 93. Subtracting the
ghosts introduces a sign problem that the old counting did not have: the ghosts are massless in
the Landau gauge, while the Goldstone modes they cancel are not, so once every mode of the
potential is heavy compared with the temperature the modes die away and the ghost term does
not, and the difference turns negative. A negative number of degrees of freedom has no
meaning: when every mode is much heavier than the temperature, the sector has left the plasma
and its entropy and energy are zero, not less than zero. The new count is therefore floored at
zero, which is the correct limit.

\paragraph{2. The Standard Model fields in the tabulated bath.} To leave room in the bath for
the Standard Model fields that the potential already contains, \code{TL} froze the tabulated
degrees of freedom above half the mass of the lightest massive \code{is\_SM} field
($0.888\,\text{GeV}$ for the 2HDM), holding them at their value there for every higher
temperature. Freezing is a blunt instrument for a selective job: it also froze the photon, the
gluons, the light fermions and the QCD crossover, none of which the potential contains. The
freezing temperature was also stored once for the whole module rather than per model, so in a
session that builds several models the last one built decided which tables all of them used:
constructing a dark model after a 2HDM removed the freezing again and silently changed the
2HDM's degrees of freedom. It is
replaced by subtracting exactly the Standard
Model fields of the potential from the tables, at the masses they have in the zero-temperature
vacuum. The inventories then close: in that vacuum the potential adds exactly what the table loses,
so the total is the full Standard Model table plus the additional fields of the model. The
subtraction depends only on the temperature once the model is built, so it is tabulated once
per model and interpolated, to a relative accuracy of $3\times10^{-7}$.

\paragraph{3. The transverse photon of the 2HDM.} The neutral gauge sector has two mass
matrices in the $(W^3, B)$ basis, not one. The longitudinal modes carry the Debye masses
$\Pi_{W_\text{L}}$ and $\Pi_{B_\text{L}}$, which are added to the matrix \emph{before} it is
diagonalised -- since the two differ, the longitudinal mixing angle depends on the temperature
and the longitudinal ``photon'' is not the photon direction. The transverse modes carry no
thermal mass at this order, so their matrix is the same one without the Debye terms. \code{TL}
took both eigenvalues of the longitudinal matrix, so the transverse photon carried the Debye
mass of the longitudinal one. The transverse matrix has entries $g_2^2\phi^2/4$,
$-g_1 g_2 \phi^2/4$ and $g_1^2\phi^2/4$; its determinant vanishes identically, so its
eigenvalues are $(g_1^2+g_2^2)\phi^2/4$ and $0$ -- the transverse $Z$, and a transverse photon
that is massless at every field value, as electromagnetism being unbroken requires. The wrong
mass entered the thermal potential and, with two degrees of freedom, the daisy term, so it
moved the barrier of every 2HDM run.

\section{The scale factor of the percolation integral}

The percolation integral needs the scale factor $a(T)$ of the false vacuum, which \code{TL}
obtained by integrating $\mathrm{d}\ln a/\mathrm{d}T = -1/(3 c_s^2 T)$ with the trapezoidal
rule over the solver's temperature grid. That grid spans the overlap of the two traced phases,
which for the 2HDM benchmark runs over thirteen decades with 71 points. Far below completion
the sound speed of the false vacuum is meaningless -- 43 of those 71 values were the silent
fallback to $1/3$ -- and the trapezoidal rule weights each interval by the value at its cold
end, so one interval alone contributed 311 e-folds: $a(T)$ jumped by $10^{135}$ between
neighbouring grid points and reached $10^{217}$, and the integrator gave up with error~10,
i.e. by an error marked (correctly) as stemming from a numerical source in the percolation solver.

The integral form of the same relation is entropy conservation, $a^3 s = \text{const}$, which
needs no integration over the temperature grid, in order to obtain the scale factor directly:
$a(T)/a(T_\text{hot}) = [s(T_\text{hot})/s(T)]^{1/3}$. The
entropy is taken from the degrees of freedom of the transitioning sector rather than from
$-\partial V/\partial T$. The reason is that the Arnold--Espinosa daisy term tends to
$-T \sum_i n_i m_i^3/(12\pi)$ once the modes are heavy, a term linear in the temperature, whose
contribution to the entropy $s = -\partial V/\partial T$ is therefore a constant. That constant
cannot be physical: a species that is much heavier than the temperature is Boltzmann
suppressed and carries an entropy that falls at least as fast as $T^3$, so a plasma of such
species cannot keep a temperature-independent entropy. It is an artefact of the
high-temperature expansion the daisy term is derived from, and for the 2HDM benchmark it
overstates the entropy of the plasma by a factor $3.7$ at $T = 1\,\text{GeV}$ and by two orders
of magnitude at $T = 0.1\,\text{GeV}$.\footnote{The artefact itself is left for the next pull
request, which switches the Debye masses off as the particles that build them leave the plasma,
following arXiv:2010.09744 and arXiv:2306.17239. Here it is only avoided, by not taking the
entropy from the potential.} Where the entropy cannot be used, the bag relation $a \propto 1/T$
continues the chain, and a step that would contract the universe as it cools is replaced by it
as well. For a pure radiation bath, where $a \propto 1/T$ is exact, the old rule was wrong by a
factor $680$ at the cold end of a twelve-decade grid; the new one is exact.

\section{What it changes}

\begin{table}[h]
\centering
\begin{tabular}{lrrrr}
\toprule
median change [\%] & counting & photon & scale factor & all three \\
\midrule
$T_\text{crit}$                & $0$     & $-0.26$  & $0$      & $-0.26$ \\
$T_\text{perc}$                & $+0.02$ & $-0.35$  & $-0.001$ & $-0.34$ \\
$T_\text{reh}$                 & $+0.04$ & $-0.27$  & $-0.001$ & $-0.27$ \\
$\alpha$                       & $+1.4$  & $+1.4$   & $+0.002$ & $+2.9$ \\
$(\beta/H)_{R_*}$              & $-1.0$  & $-0.2$   & $+0.03$  & $-0.6$ \\
$g_\text{eff}$ after reheating & $-2.9$  & $-0.05$  & $-0.000$ & $-2.9$ \\
$h_\text{eff}$ after reheating & $-1.8$  & $-0.05$  & $-0.000$ & $-1.8$ \\
peak amplitude                 & $+3.9$  & $+3.7$   & $-0.05$  & $+7.1$ \\
\bottomrule
\end{tabular}
\caption{2HDM type~I line of the code comparison, 49 values of $\lambda_3$, against the state
before this pull request. The first column holds both halves of the counting, the modes of the
potential and the subtraction from the tabulated bath, which were validated together; the
second the transverse photon; the third the scale factor, which changes nothing but gives back
the one point that the photon fix costs, so the line has 49 of 49 results.}
\end{table}

In the dark sector models implemented in the code, the percolation is essentially untouched:
along the abelian dark Higgs line $T_\text{perc}$, $T_\text{reh}$, $\alpha$ and $\beta/H$ are
bit-identical, and along the conformal line $T_\text{perc}$ moves by $10^{-5}\,\%$ and
$T_\text{reh}$ by $+0.016\,\%$. What does move is the redshift: the degrees of freedom after
reheating change by $-0.9\,\%$ (conformal) and $+0.3\,\%$ (abelian dark Higgs). Where the Goldstone modes are massless at the minimum the counting is unchanged to
machine precision; where they are not, as in the conformal dark U(1) model, it changes by
$16\,\%$ at $T = 0.1\,v$.

\paragraph{Still open.} The daisy term is derived in the high-temperature expansion and leaves
spurious temperature-independent pieces behind at $T \ll m$, which is why the sound speed
of the broken phase drifts away from $1/3$ at low temperature instead of towards it. A
following pull request will switch the Debye masses off as the particles that build them leave
the plasma (see arXiv:2010.09744 and arXiv:2306.17239), validated on its own.

\end{document}
